\section{Statement of the Problem}

Inconsistent feeding and drinking practices in pigs can lead to significant health risks and feed wastage \cite{oecdfao2025outlook}. Early detection of illness through physical or manual observation is often unreliable because subtle symptoms may go unnoticed \cite{usda2026pork}. Without a consistent tracking system, caretakers must rely on labor-dependent observation, which can result in inaccurate records and delays in identifying symptoms or other problems. A pig's health status is closely linked to its feeding and drinking behavior, as measurable deviations in dietary intake are often among the first observable indicators of illness \cite{dingcong2017monitoring}. However, without a data-recording system, individual feeding patterns are not established or tracked over time, making it difficult to distinguish normal behavioral variation from significant decline. In addition, unmonitored air quality in piggery environments poses an unnecessary and preventable risk of common respiratory illnesses such as coughs and colds among pigs \cite{livestockproduction2026}.

This study addresses these issues by proposing a system that establishes per-module consumption baselines using a rolling time window and flags deviations beyond a defined threshold as potential signs of abnormal behavior requiring immediate caretaker attention. Air-quality readings from integrated sensors are also compared against safe-range standards, with alerts generated when toxicity levels within the pig-pen environment exceed acceptable limits. This approach aims to reduce reliance on subjective visual observation and provide a more consistent and reliable basis for early health monitoring.
